\section{Compilation into one Boolean oracle}\label{sec:compilation}

The skeleton of Proposition~\ref{rec:skeleton} specifies only polynomially
many slots, although each slot may have a large table. We now encode finite
gate words for those tables in one Boolean function. Two details require
care: the words must approximate matrices with their actual phases, and the
resulting circuit must be generated without finding those words. The need to
retain scalar phases under coherent control is already present in the general
controlled-gate decompositions of~\cite{Barenco1995}.

\subsection{Phase-sensitive one-qubit approximation}

We first prove the approximation fact needed for the compilation. The
argument combines the irrational-rotation proof of universality for this
gate set~\cite{Boykin2000} with the commutator refinement underlying the
Solovay--Kitaev theorem~\cite{DN2006}. The weaker polynomial bound below is
sufficient for our purpose.

\begin{lemma}\label{compile:words}
There is an absolute positive integer $K_{\mathrm w}$ such that every $S\in SU(2)$ and every
$0<\xi<1$ admit a word in $H,T,T^\dagger$ of length at most $K_{\mathrm w}\xi^{-3}$
whose matrix has determinant one and differs from $S$ by at most $\xi$ in
operator norm.
\end{lemma}

\begin{proof}
Write $X_p,Y_p,Z_p$ for the Pauli matrices, with $Y_p=iX_pZ_p$, and put
$\boldsymbol\sigma=(X_p,Y_p,Z_p)$. We first show that the determinant-one
words are dense in $SU(2)$. Up to their scalar factors,
\[
 T=e^{i\pi/8}e^{-i(\pi/8)Z_p},\qquad
 HTH=e^{i\pi/8}e^{-i(\pi/8)X_p}.
\]
Thus $W=(HTH)T=e^{i\pi/4}V$, where $V\in SU(2)$ has trace
\[
 \Tr V=2\cos^2(\pi/8)=1+\frac{\sqrt2}{2}.
\]
Recall that an algebraic integer is a root of a monic polynomial with
integer coefficients; equivalently, its monic minimal polynomial over
$\mathbb Q$ has integer coefficients. The displayed trace is not an algebraic
integer, because its monic irreducible polynomial is $x^2-2x+1/2$.
Consequently the eigenvalues of $V$ cannot be roots of unity: roots of unity
are algebraic integers, and sums of algebraic integers are algebraic integers.
Writing its eigenvalues as $e^{\pm i\theta}$, we have
$\theta/\pi\notin\mathbb Q$. The word $W^4=-V^4$ has determinant one,
and its powers are dense in the one-parameter circle about the rotation
axis of $V$.

For completeness, every matrix in $SU(2)$ has the form
$\cos t\,I+i\sin t\,\mathbf v\cdot\boldsymbol\sigma$, with
$\mathbf v\in\R^3$ a unit vector, except that the axis can be chosen
arbitrarily when $\sin t=0$. Multiplying the two rotations defining $V$
shows that its axis is not parallel to the $Z_p$ axis. Conjugation by $T$
rotates this axis through $\pi/4$ about $Z_p$, giving a nonparallel axis.
The closure of the determinant-one words contains both corresponding
circles: conjugation of a determinant-one word by $T$ is again such a word.
Rotating the second axis slightly about the first moves it out of the
plane they span. Conjugation by an element of the first circle therefore
gives a third circle with a linearly independent axis.
Let $\mathbf v_1,\mathbf v_2,\mathbf v_3$ be the three axes. The derivative
at zero of
\[
 (t_1,t_2,t_3)\longmapsto
 \prod_{j=1}^3\exp(i t_j\mathbf v_j\cdot\boldsymbol\sigma)
\]
sends $(t_1,t_2,t_3)$ to
$i(\sum_jt_j\mathbf v_j)\cdot\boldsymbol\sigma$, an isomorphism onto the
three-dimensional tangent space of $SU(2)$ at the identity. The inverse
function theorem shows that these products contain a neighborhood of the
identity. The closure of our words is thus an open subgroup. Every other
coset is open as well, so connectedness of $SU(2)$ forces this subgroup to
be all of $SU(2)$. Here connectedness also follows from the preceding
representation of $SU(2)$ as the unit sphere in $\R^4$.

We next obtain a uniform length bound. Density and compactness provide a
finite $\varepsilon_0$-net of determinant-one words of some fixed maximum
length $\ell_0$, where $\varepsilon_0>0$ will be chosen sufficiently small.
Suppose that words of length at most $\ell$ form an $\varepsilon$-net, and
choose such a word $W_0$ within $\varepsilon$ of a target $S$. The residual
$S W_0^\dagger$ has the form
\[
 S W_0^\dagger=\exp(i\mathbf a\cdot\boldsymbol\sigma),
 \qquad |\mathbf a|\le 2\varepsilon,
\]
for sufficiently small $\varepsilon$. The Pauli commutator identity gives
skew-Hermitian matrices
\[
 P=i\mathbf u\cdot\boldsymbol\sigma,\qquad
 Q=i\mathbf v\cdot\boldsymbol\sigma,
 \qquad [P,Q]=i\mathbf a\cdot\boldsymbol\sigma,
\]
with $\norm{P},\norm{Q}=O(\sqrt\varepsilon)$: take perpendicular vectors
with $\mathbf u\times\mathbf v=-\mathbf a/2$ and equal lengths.
For $A=e^P$ and $B=e^Q$, expansion of the four exponential series through
degree two gives
\[
 ABA^\dagger B^\dagger=I+[P,Q]+O(\varepsilon^{3/2})
                    =S W_0^\dagger+O(\varepsilon^{3/2}).
\]
All errors here are in operator norm with absolute constants; the remaining
terms of the exponential series are bounded by a constant times
$(\norm P+\norm Q)^3$.

Choose net words $\widetilde A,\widetilde B$ within $\varepsilon$ of
$A,B$. The gain in precision survives this substitution. Indeed,
\begin{align*}
 \norm{ABA^\dagger B^\dagger-
        \widetilde A B\widetilde A^\dagger B^\dagger}
 &\le 2\varepsilon\norm{B-I},\\
 \norm{\widetilde A B\widetilde A^\dagger B^\dagger-
        \widetilde A\widetilde B\widetilde A^\dagger\widetilde B^\dagger}
 &\le 2\varepsilon\norm{\widetilde A-I}.
\end{align*}
The first inequality follows by writing $B=I+(B-I)$ and cancelling the
identity terms; the second is the same calculation with the roles reversed.
Both right sides are $O(\varepsilon^{3/2})$. Therefore the word
\[
 \widetilde A\widetilde B\widetilde A^\dagger
 \widetilde B^\dagger W_0
\]
approximates $S$ to $C\varepsilon^{3/2}$ and has length at most $5\ell$.
Inverse words are available because $H^\dagger=H$ and both $T$ and
$T^\dagger$ belong to the alphabet.

Choose the fixed initial net so small that
$C\sqrt{\varepsilon_0}\le1/2$. Repeating the refinement $j$ times gives
error at most $2^{-j}\varepsilon_0$ and length at most $5^j\ell_0$.
Taking $j$ just large enough for the former quantity to be at most $\xi$
gives a length bounded by $K_{\mathrm w}\xi^{-3}$, since $\log_2 5<3$.
\end{proof}

Only the existence of the fixed integer $K_{\mathrm w}$ will be used. A generator can
hardcode this integer; it does not need the net or a procedure
that selects approximating words. Those selections will be stored in the
oracle.

\subsection{Lookup, control, and restoration of scratch space}

\begin{proposition}[Compilation of slots]\label{compile:slots}
Let a skeleton have $S$ slots on $Q$ qubits, and put $\eta=1/h$ for an integer
$h\ge2$.
There is a deterministic procedure, polynomial in $Q,S,1/\eta$, that produces
a circuit over $H,T,T^\dagger,\mathrm{CNOT}$ and one Boolean oracle, with
polynomially many qubits, gates, calls, and address bits. For every assignment
of the slot tables, some contents of that oracle make the circuit differ
from the prescribed skeleton by at most $S\eta$ on inputs with its added
scratch qubits zero, with the comparison circuit keeping those scratch
qubits zero.
\end{proposition}

\begin{proof}
First consider one table entry $W_y\in U(2)$. Choose a real number $\gamma_y$
and a matrix $S_y\in SU(2)$ such that
\[
 W_y=e^{i\gamma_y}S_y.
\]
By Lemma~\ref{compile:words}, choose words within $\eta/2$ of $S_y$ and of
\[
 P_y=\diag(e^{i\gamma_y},e^{-i\gamma_y}).
\]
Pad every word with identities to the common length
$\ell=\lceil K_{\mathrm w}(2/\eta)^3\rceil$. Apply the first word to the slot target
and the second to a fresh scratch qubit initially in $\ket0$. Their ideal
actions on this input multiply to $W_y$ on the target and return scratch
to zero. The two approximation errors add to at most $\eta$. Since
different $y$ occupy orthogonal control blocks, this is an operator-norm
bound for the entire controlled slot. In particular, the scalar factor
$e^{i\gamma_y}$ has been retained as a relative phase between control
values.

Use two bits to encode each symbol in $\{I,H,T,T^\dagger\}$. A single
Boolean function stores all symbol bits at addresses containing
\[
 (\text{slot number},\ y,\ \text{which word},\
   \text{word position},\ \text{symbol bit}).
\]
Pad the field for $y$ to $Q$ bits. This makes every address have a common
length $O(Q+\log(S+1)+\log(\ell+1))$. For each word position, copy the
slot's controls into the address by CNOTs, set the fixed address bits, and
query the two symbol bits into zero qubits. Apply the specified symbol
conditionally, then erase the symbol bits by querying again and reverse
the address preparation. The original slot controls do not change under
any of these target operations, so the erasure is exact even in a
superposition of control values.

We verify that symbol selection uses only the allowed elementary gates.
Bit flips are available as $HT^4H=X_p$. On bits $a,b,c$, the integer identity
\begin{equation}\label{compile:parity}
 4abc=a+b+c-(a\oplus b)-(a\oplus c)-(b\oplus c)
                         +(a\oplus b\oplus c)
\end{equation}
implements a doubly controlled $Z_p$ using $T,T^\dagger$: compute each
parity into clean scratch, apply its indicated phase, and uncompute it.
Conjugating the target by $H$ gives Toffoli. Thus equality tests on symbol
bits and all required reversible address operations are exact.
To apply $T$ conditionally, compute the AND of its control and target into
a clean scratch bit, apply $T$ there, and uncompute the AND; the same
construction applies to $T^\dagger$.

For a controlled $H$, define
\[
 D_0=(T^2H)T(T^2H)^\dagger.
\]
Because $(T^2H)Z_p(T^2H)^\dagger=Y_p$, we have
$D_0=e^{i\pi/8}e^{-i(\pi/8)Y_p}$ and hence
\[
 D_0Z_pD_0^\dagger=(X_p+Z_p)/\sqrt2=H.
\]
Apply $D_0^\dagger$ to the target, then a controlled $Z_p$, then $D_0$.
When the control is zero the two unconditional gates cancel; when it is
one their product is exactly $H$. A controlled $Z_p$ itself is a CNOT
conjugated on its target by $H$. This proves that every selected symbol,
including its phase, can be applied exactly.

It remains to compose the approximate slots. Let $A_j$ denote the actual
implementation of slot $j$, and let $V_j$ denote the ideal slot acting
trivially on all added scratch. The phase scratch may be reused. Although
its actual state need not be exactly zero after a slot, the identity
\[
 A_S\cdots A_1-V_S\cdots V_1
 =\sum_{j=1}^S A_S\cdots A_{j+1}(A_j-V_j)V_{j-1}\cdots V_1
\]
uses the error $A_j-V_j$ only after an ideal prefix. Such a prefix keeps
all added scratch at zero, precisely where the one-slot bound applies.
The unitary suffix preserves its norm. The total error is therefore at
most $S\eta$ on the claimed input subspace.

There are two words per slot, $O(\ell)$ symbol positions per word, and a
polynomial number of elementary operations per lookup. The addresses,
scratch, and total circuit size are thus polynomial in $Q,S,1/\eta$.
Generating the circuit requires only the skeleton, the fixed length bound,
and ordinary integer arithmetic. In particular, it never selects a word
or computes a table entry. Each target merely determines one admissible
Boolean function storing all choices. Adjoint slots can have their own
addresses and approximating words, or can reverse and invert stored
words; either convention uses the same single oracle.
\end{proof}

\begin{proof}[Proof of Theorem~\ref{thm:main}]
Apply Proposition~\ref{rec:skeleton} with $d=2^n$ and no initial labels.
Its skeleton has polynomially many qubits and slots, and approximates the
target unitary, including restoration of its working qubits, to error
less than $1/100$. If it has $S$ slots, compile it using
\[
 \eta=\frac{1}{100(S+1)}.
\]
Proposition~\ref{compile:slots} adds error less than $1/100$. The entire
circuit is therefore within $1/10$ in operator norm, on zero-workspace
inputs, of applying $U$ and returning every other qubit to zero. The same
polynomial bounds its width, number of gates and calls, oracle address
length, and the running time of its deterministic generator, after
increasing fixed coefficients if necessary.

Take the original index qubits as the designated outputs. For a pure input
state $\ket\psi$ together with a reference register, let $a$ and $b$ be the
actual and ideal final unit vectors before discarding workspace. The
operator-norm estimate gives $\norm{a-b}\le1/10$. Since the trace norm of
a rank-one matrix $uv^\dagger$ is $\norm u\norm v$,
\[
 \norm{aa^\dagger-bb^\dagger}_1
 \le\norm{(a-b)a^\dagger}_1+\norm{b(a-b)^\dagger}_1
 \le2\norm{a-b}\le\frac15.
\]
Partial trace cannot increase the trace norm of a Hermitian matrix, and
convexity gives the same bound for mixed input states. In particular the
supremum in Theorem~\ref{thm:main}, with its $n$-qubit reference, is at most
$1/5<1/2$.
\end{proof}
